\documentclass[11pt]{article}
\newcommand{\runningtitle}{Geometric progressions over finite fields}
\usepackage[T1]{fontenc}
\usepackage{lmodern,amsmath,amsthm,amssymb,booktabs,microtype,xcolor}
\usepackage[letterpaper,margin=.94in,top=.83in,bottom=.86in]{geometry}
\usepackage[colorlinks=true,linkcolor=black,urlcolor=teal!65!black,citecolor=black,breaklinks=true]{hyperref}
\usepackage{fancyhdr}
\pagestyle{fancy}
\fancyhf{}\lhead{\small\sffamily\runningtitle}\rhead{\small\sffamily Research note}\cfoot{\small\thepage}
\setlength{\headheight}{14pt}
\setlength{\emergencystretch}{3em}
\newtheorem{theorem}{Theorem}
\newtheorem{lemma}{Lemma}
\newcommand{\paperdate}{November 25, 2026}
\date{\paperdate}
\author{}
\raggedbottom

\hypersetup{pdftitle={A counterexample to the geometric-progression density conjecture over finite fields}}
\title{A counterexample to the geometric-progression density conjecture over finite fields}
\begin{document}
\maketitle
\begin{abstract}
For every prime power $q$, we construct a subset of $\mathbb F_q[x]$ free of three-term geometric progressions with nonconstant polynomial ratio whose upper density exceeds the value proposed by Asada, Fourakis, Manski, McNew, Miller and Moreland. The proof combines a carry-free sphere of digit vectors with a change in the allowed degrees for polynomials having a specified rough part. A generating function converts an additive gain into a strict density gain. We also give finite optimization problems converging to the true supremum with an explicit error bound.
\end{abstract}

\section{A counterexample to the polynomial-ring density conjecture}
\label{sec:density}
For every prime power $q$, there is a subset of $\mathbb F_q[x]$ containing
no triple $f,rf,r^2 f$ with $f\ne0$ and $r$ a nonconstant polynomial, whose
upper density strictly exceeds
\[
 m_q=(1-q^{-1})\sum_{a\in A}q^{-a}
     =(1-q^{-2})\prod_{j\ge1}(1+q^{-3^j}),
 \qquad A=\{a\ge0:\text{all ternary digits of }a\text{ are }0\text{ or }1\}.
\]
Here upper density means $\limsup_{N\to\infty}|E\cap\{f:\deg f\le N\}|/q^{N+1}$.
This disproves Conjecture 6.3 of Asada, Fourakis, Manski, McNew, Miller and
Moreland, \emph{Subsets of $\mathbb F_q[x]$ free of 3-term geometric progressions},
\emph{Finite Fields Appl.} 44 (2017), 135--147,
\href{https://doi.org/10.1016/j.ffa.2016.10.002}{doi:10.1016/j.ffa.2016.10.002}.
Their Theorem 4.3 proves that $m_q$ is optimal among sets defined solely by degree.
Our construction uses the factorization of polynomials of the same degree.
An earlier ProveIt analysis, \texttt{progressions.tex} inside
\texttt{grh-progression-densities-and-cubic-prime-race.zip}, improved other bounds
and explicitly left Conjecture 6.3 open. Its use of exponents suggested revisiting
the fixed-radius digit construction in \texttt{exists\_three\_term\_free}
in the readable Lean edition. The construction there selects embeddings of CM fields.
Its elementary combinatorial ingredient, the classical \href{https://doi.org/10.1073/pnas.32.12.331}{Behrend sphere method},
here selects degrees inside a class of smooth polynomials.

\subsection*{An arithmetic-progression-free replacement}
In a ball of cutoff degree $N$, call $d=N-\deg f$ the deficit.
The rule $d\in A$ gives a progression-free set with density tending to $m_q$ as
$N\to\infty$. Indeed, nonconstant polynomial ratios give strictly increasing
degrees in arithmetic progression, and the equation $a+c=2b$ for $a,b,c\in A$
has no ternary carries, forcing equality at every digit.

Put $M=3^{20}$. For $u=(u_0,\ldots,u_8)\in\{0,\ldots,5\}^9$, write
$b(u)=\sum_{i=0}^8u_i11^i$. Take the set $B$ of these integers satisfying
$\sum_i(2u_i-5)^2=105$. The one-digit squared contributions are $1,9,25$,
each with multiplicity two. If their respective counts are $i,j,k$, then
$i+j+k=9$ and $j+3k=12$. The possibilities $(i,j,k)$ are $(1,6,2)$,
$(3,3,3)$ and $(5,0,4)$. Hence
\[
 |B|=2^9\!\left(\frac{9!}{1!6!2!}+\frac{9!}{3!3!3!}
                       +\frac{9!}{5!0!4!}\right)
     =1{,}053{,}696 > 2^{20}=|A\cap[0,M)|.
\]
Also $\max B\le(11^9-1)/2<M/2$. Neither the sum of two digits nor a
doubled digit exceeds 10. Thus a midpoint identity implies a
coordinatewise midpoint identity. Strict convexity of the squared norm
on a line forces all three centered vectors to coincide. This is the
part of the proof shared with the source Lean construction.

We claim that $D=B\cup(A\cap[M,\infty))$ is also progression-free.
Two low terms cannot reach a high term, since twice the largest low term
is less than $M$. Suppose $a<b<c$ with $a\in B$ and $b,c\in A$ high.
Write $S=3^j$ for the largest ternary place of $c$, so that
$S\le c\le(3S-1)/2$. If $b<S$ then $2b\le S-1<c$.
If $b\ge S$ then $2b-c\ge(S+1)/2>M/2>a$.
Neither case admits $a+c=2b$. Triples lying entirely on one side
are excluded by the progression-free property already proved.

\subsection*{A smooth factor of fixed degree records a strict gain}
Let $\mathcal P$ be the finite set of all monic irreducibles in
$\mathbb F_q[x]$ of degree at most $M$. Call a monic polynomial smooth
if all its factors belong to $\mathcal P$ and rough if none do.
Let $H_n$ count the monic smooth polynomials of degree $n$, with
$H_n=0$ for $n<0$ and $H_0=1$. Unique factorization gives
\[
 H(z)=\sum_{n\ge0}H_n z^n=\prod_{P\in\mathcal P}(1-z^{\deg P})^{-1}.
\]
For $0<z<1$ this series is absolutely convergent. The polynomial
$J(z)=\sum_{d\in B}z^d-\sum_{d\in A,\,d<M}z^d$ satisfies $J(1)>0$.
For real $z<1$ sufficiently near 1 we have $H(z)J(z)>0$. Its coefficients
cannot all be nonpositive. Fix an $R\ge0$ for which
\[
 \Delta=[z^R]H(z)J(z)
   =\sum_{d\in B}H_{R-d}-\sum_{d\in A,\,d<M}H_{R-d}\ \ge1.
\]
All parameters so far depend at most on $q$. They are finite and can in
principle be found by integer enumeration. The positivity argument avoids
evaluating an infeasibly large coefficient.

An explicit bound on the search is
$R\le R_* = M(1+2^{20}|\mathcal P|)$. To see this, put
$z=1-1/(1024M)$ and give $T=n$ probability $H_nz^n/H(z)$.
Logarithmic differentiation at $z$ gives
\[
 \mathbb E T=\sum_{P\in\mathcal P}
 \frac{(\deg P)z^{\deg P}}{1-z^{\deg P}}
 \le \frac{|\mathcal P|}{1-z}=1024M|\mathcal P|.
\]
Indeed $1-z^e=(1-z)\sum_{j<e}z^j\ge e(1-z)z^e$ for $e\ge1$.
By Markov's inequality $\Pr(T\le R_*-M)\ge1023/1024$.
Also $z^M\ge1023/1024$ by Bernoulli's inequality. In the sum of
coefficients of $H J$ through $R_*$ weighted by $z^n/H(z)$, each added
term $d\in B$ contributes $z^d\Pr(T\le R_*-d)\ge(1023/1024)^2$.
Each displaced term contributes at most one in absolute value. The
weighted sum is strictly positive since $1023^2|B|>1024^2\,2^{20}$.
This proves the asserted bound.

At cutoff $N\ge R$ start with the deficit rule $A$. For every squarefree rough
monic $g$ of degree $N-R$, consider all $cgh$ in the ball, where $c\in\mathbb F_q^\times$
and $h$ is monic smooth. On each such class replace the rule by $d\in D$.
Every nonzero polynomial has a unique scalar and a unique rough and smooth
part, so the classes are disjoint. For each $g$ the change is exactly
$(q-1)\Delta$ polynomials, since a term of deficit $d$ has $\deg h=R-d$.

To verify avoidance, any newly created progression would contain an added
polynomial $f=cgh$. Its deficit belongs to $B$ and is less than $M/2$.
If $f$ is first or middle, the next term has degree at most $N$, so
$\deg r\le N-\deg f<M$. If $f$ is last, then $r^2$ divides $f$.
Every irreducible outside $\mathcal P$ occurs in $f$ with exponent at
most one, by squarefreeness of $g$. Consequently all nonconstant factors
of $r$ belong to $\mathcal P$ in every case. All three terms have the
same monic rough part $g$. Their deficits would be three distinct members
of $D$ in arithmetic progression, a contradiction. This argument uses the
condition $r\in\mathbb F_q[x]$ of Conjecture 6.3.

\subsection*{Positive density and passage to one infinite set}
If $G_n$ counts squarefree rough monic polynomials of degree $n$, then
\[
 \frac{G_n}{q^n}\longrightarrow
 c_q=(1-q^{-1})\prod_{P\in\mathcal P}(1+q^{-\deg P})^{-1}>0 .
\]
Indeed their generating function is
$(1-qz^2)(1-qz)^{-1}\prod_{P\in\mathcal P}(1+z^{\deg P})^{-1}$.
This follows from the Euler product $\prod_P(1-z^{\deg P})^{-1}=(1-qz)^{-1}$
and $1+y=(1-y^2)/(1-y)$. Removing the pole at $1/q$ leaves a rational
function analytic on $|z|<1$, giving the displayed limit by coefficient
comparison. If $E_N$ is the resulting finite set, division by the ball size gives
\[
 \frac{|E_N|}{q^{N+1}}
 =(1-q^{-1})\sum_{\substack{d\in A\\d\le N}} q^{-d}
     + \frac{(q-1)\Delta G_{N-R}}{q^{N+1}}
 \longrightarrow m_q+\underbrace{(q-1)\Delta q^{-R-1}c_q}_{\displaystyle\delta_q>0}.
\]
This establishes a fixed positive margin for each $q$. In particular the gap is at least $(q-1)q^{-R_*-1}c_q$, which involves only explicitly specified finite products and integers.

Choose integer cutoffs $N_j$ increasing sufficiently fast that
$N_j-j>2N_{j-1}$ and $N_j-2j>N_{j-1}$, with $N_1\ge R+2$.
Take from $E_{N_j}$ only polynomials with degrees between $N_j-j$ and
$N_j$ inclusive, and form the union of these blocks.
A triple meeting an earlier block and the latest block in two earlier
terms would have final degree at most $2N_{j-1}$. With two latest terms,
its first degree would be at least $N_j-2j$. Both are impossible.
Deleting the degrees below $N_j-j$ loses at most $q^{N_j-j}$ polynomials at
that cutoff. The normalized loss is at most $q^{-j-1}$, which tends to zero.
The union has upper density at least $m_q+\delta_q$, as claimed.

The proof gives a positive gap with a very large cutoff; it leaves the true
supremum undetermined. The small digit computation in \texttt{check\_sphere.py}
independently verifies the cardinality and the range used in the replacement.

\subsection*{A smaller certified cutoff}
The same construction admits $M=3^8=6561$. Take the ordinary three-term-progression-free sets listed
as $a(106)\le795$ and $a(151)\le1241$ in
\href{https://www.math.uni.wroc.pl/~jwr/non-ave/DATABASE.TXT}{Wróblewski's database of nonaveraging sets}.
Translate each to start at zero and call them $U$ and $V$, respectively.
The accompanying \texttt{check\_small\_pattern.py} contains both sets in full,
encoded as integer bitmasks, and checks every pair of possible first and middle terms.
Their sizes are 106 and 151 and their respective maxima are 794 and 1240. Set
\[
 B=U\ \cup\ (2035+V).
\]
The distance between the two blocks is $2035-794=1241$.
Two elements in one block differ by at most 1240, so a nonconstant arithmetic
progression of length three cannot meet both blocks. Consequently $B$ is
progression-free, $|B|=257>2^8$, and $\max B=3275<M/2$.
These finite checks are sufficient for the replacement argument. They use only
the displayed entries of the database, independently of any optimality claims.

The smaller cutoff also gives the finite search bound
\[
 R\le M\bigl(1+257\cdot514\,|\mathcal P|\bigr),
 \qquad \mathcal P=\{P:\ P\text{ monic irreducible},\ \deg P\le6561\}.
\]
Indeed, take $z=1-1/(257M)$ in the probability distribution on $T$ above.
Then $\mathbb E T\le257M|\mathcal P|$, so $\Pr(T\le257\cdot514M|\mathcal P|)
\ge513/514$. For each added $d<M/2$, Bernoulli's inequality gives
$z^d\ge1-d/(257M)\ge513/514$. The truncated weighted coefficient sum is
positive because $257\cdot513^2-256\cdot514^2=257>0$.
This supplies a version of the explicit density gap using the smaller finite
set of irreducibles.

\section{A weighted replacement using irreducibles of degrees at most 480}
\label{sec:tilted}
A finite arithmetic-progression-free set $B$ of nonnegative deficits is
enough whenever there is a real $0<z<1$ with
\begin{equation}
 J(z)=\sum_{b\in B}z^b-\sum_{a\in A}z^a>0 .
 \label{eq:tilt}
\end{equation}
Indeed take $D=B$ on the modified classes and keep the original rule on
the remaining polynomials. Use the smooth primes of degree at most $\max B$.
At fixed smooth cutoff $R$, the class gain per scalar is still
$[u^R]H(u)J(u)$. Positive evaluation at $u=z$ implies that some coefficient
is positive. The progression argument needs only the upper bound on newly
added deficits and the fact that $B$ is progression-free. It therefore
applies without joining $B$ to a ternary tail.

Here is a witness involving 74 deficits. In J. Wr\'oblewski's
\href{https://www.math.uni.wroc.pl/~jwr/non-ave/DATABASE.TXT}{database of
non-averaging sets}, take Gavin Theobald's entry
\texttt{a(143)<=1360}, marked \texttt{1092044464}. Keep its entries
between 878 and 1358 inclusive and subtract 878. The resulting $B$ has
minimum zero and maximum 480. The bundled \texttt{check\_tilted.py}
includes the source slice, checks every midpoint, and verifies with
integer arithmetic that for $z=127/128$,
\[
 \sum_{b\in B}z^b-\prod_{j=0}^6(1+z^{3^j})>0.016,
 \qquad 128z^{2187}<0.00001 .
\]
The product counts the terms of $A$ below $3^7$. The sum of all remaining
powers is at most $z^{2187}/(1-z)$. Thus $J(z)>0.015$.

Put $s=|\mathcal P|$ for the number of monic irreducibles of degree at most
480. One may take $R\le R_0=480+2{,}000{,}000s$. To see this, assign degree
$n$ probability $H_nz^n/H(z)$ and call the resulting random degree $X$.
By the product formula the $s$ multiplicities are independent geometric
variables. A prime of degree $k$ contributes $kz^k/(1-z^k)\le z/(1-z)=127$
to $\mathbb E X$, since $1+z+\cdots+z^{k-1}\ge kz^{k-1}$.
In the normalized sum $H(z)J(z)/H(z)$, the positive part discarded by
truncating at total degree $R_0$ is bounded above by
\[
 \sum_{b\in B}z^b\Pr(X>R_0-b)
 \ \le\ 74\frac{127s}{2{,}000{,}000s}<0.005 .
\]
Discarding negative contributions can only raise the truncated sum.
Consequently its value is positive and some integral coefficient through
$R_0$ is at least one. In the lower bound $(q-1)c_q q^{-R_0-1}$ for the
density improvement, $c_q$ now uses only irreducibles of degree at most 480.

\subsection*{Why the exchange exists}
The weighted criterion also explains the failure of the conjecture without a
choice of particular digit counts. More generally let $A_0$ be any
three-term-arithmetic-progression-free subset of the nonnegative integers
with $|A_0\cap[0,x]|\ll (1+x)^\alpha$ for some $0\le\alpha<1$.
Its deficit rule has limiting ball density
$(1-q^{-1})\sum_{d\in A_0}q^{-d}$. For every $q$ one can construct a
polynomial-ratio progression-free set of strictly greater upper density.

Indeed, for $X\ge1$, Abel summation gives
$\sum_{d\in A_0}e^{-d/X}\ll_\alpha X^\alpha$, with a constant also
depending on the counting bound. Choose an integer $h\ge2$ so large that
$h>(2h-1)^\alpha$. In dimension $n$ read the digits $0,\ldots,h-1$
in base $2h-1$ and set $X=(2h-1)^n$. One squared-norm level contains at
least $h^n/(1+n(h-1)^2)$ vectors. Its integer encodings form a
three-term-free set $B$ in $[0,X/2]$ by the same carry-free midpoint
argument as before. Consequently
\[
 \sum_{d\in B}e^{-d/X}\ge
 \frac{e^{-1/2}h^n}{1+n(h-1)^2}>\sum_{d\in A_0}e^{-d/X}
\]
for all sufficiently large $n$. Use $z=e^{-1/X}$ in the positive-coefficient
argument for the weighted exchange, with $A_0$ in place of $A$. The proof
of confinement only uses the finiteness of $B$ and progression avoidance
inside and outside the modified classes. The same separated degree blocks
give the infinite set.

For the ternary $A$ in the conjecture one can take $\alpha=\log_3 2$,
because there are $2^k$ permitted integers below $3^k$. Already $h=4$
satisfies $4>7^{\log_3 2}$. Thus the transferable ingredient is a growth
gap between two explicit digit constructions. The numerical certificates
give finite versions of this argument.

\section{Computing the extremal upper density}
There is an integer precedent for computing extremal upper densities by
optimizing over a finite set of primes and then enlarging it.
See N. McNew, \emph{On sets of integers which contain no three terms in
geometric progression}, Math. Comp. 84 (2015), 2893--2910,
\href{https://arxiv.org/abs/1310.2277}{Theorems 4.2 and 5.1}.
The earlier incoming \texttt{triangle.tex} uses one and two linear primes
to bound the supremum over polynomials. The following adaptation to all
irreducibles of degree at most $T$ gives a bound on the approximation error
in terms of the degree cutoff.
Fix $q$. Write $\beta$ for the supremum in Conjecture 6.3.
For an integer $T\ge1$, let $S$ comprise all monic irreducibles of degree at
most $T$, and let $J_S(t)$ be the largest size of a progression-free subset
of the monic $S$-smooth polynomials of degree at most $t$. Put
\[
 U_T=(1-q^{-1})\prod_{P\in S}(1-q^{-\deg P})
              \sum_{t\ge0}J_S(t)q^{-t},\qquad
 v_T=\frac{1-q^{-1}}{\prod_{P\in S}(1-q^{-2\deg P})}.
\]
Then
\[
 v_T U_T\le\beta\le U_T,\qquad
 0\le U_T-\beta\le 1-v_T\le \frac{q^{-T}}{q-1}.
\]
Thus $\beta=\lim_{T\to\infty}U_T$ is a computable real. The formula
also holds with rational-function ratios if that convention is used for
both $\beta$ and every $J_S(t)$.

For the upper bound, fix the scalar and the monic rough part $g$.
At cutoff $N$, at most $J_S(N-\deg g)$ smooth factors can occur in that
class. For the lower bound, restrict the allowed $g$ to be squarefree
and use a maximizing subset independently in each class, using the same
choice for all scalars at fixed $g$. In any allowed progression the valuations
outside $S$ are three equally spaced members of $\{0,1\}$, hence constant.
The three monic parts therefore lie in one class, where their quotients
would give the prohibited progression between smooth factors. This also
handles signed valuation steps in the rational-ratio version, including
nonconstant ratios of degree zero. To glue the finite constructions,
use the degree blocks from Section~\ref{sec:density}. A rational-ratio
progression with nonconstant degrees can be ordered in increasing degree
by reversing its ratio if necessary. Equal-degree triples are already
confined to one block.

The proportion of monic degree-$n$ polynomials coprime to all primes in
$S$ tends to $c=\prod_{P\in S}(1-q^{-\deg P})$, by finite
inclusion--exclusion. The proportion also squarefree tends to $cv_T$ by
the generating function in Section~\ref{sec:density}. In the normalized
count, the classes with $\deg g=N-t$ therefore contribute at most
$(1-q^{-1})c\,J_S(t)q^{-t}$ in the limit, with the corresponding
lower construction attaining the same expression times $v_T$. Dominated
convergence applies since $J_S(t)\le(t+1)^{|S|}$ by counting the possible
exponent vectors. Also $U_T\le1$, by applying the identical count with
all smooth factors allowed. Finally
\[
 1-v_T=1-\prod_{P\notin S}(1-q^{-2\deg P})
 \le \sum_{d>T}q^d q^{-2d}=\frac{q^{-T}}{q-1}.
\]
For each fixed $T$, finite enumeration of admissible sets computes
$J_S(t)$ exactly. Truncating its series at $t=L$ costs at most the
explicit tail $\sum_{t>L}(t+1)^{|S|}q^{-t}$, which tends effectively to zero.
This characterizes the supremum to arbitrary prescribed error. Obtaining
a practically sharper numerical interval for small $q$ requires solving
large finite extremal problems.

The supremum is attained as an upper density. For stage $j$ take $T=j$,
and choose $N_j$ sufficiently large that the squarefree construction has
density at least $v_jU_j-1/j$ in its ball and satisfies both block-separation
inequalities of Section~\ref{sec:density}. Keep its top $j+1$ degree layers.
Their union is progression-free and has upper density at least
$\lim_{j\to\infty}(v_jU_j-1/j-q^{-j-1})=\beta$. By definition of
$\beta$ equality follows.

\section{How large can the improvement be?}
For $q=2$, let $\beta_2$ denote the true supremum. The earlier two-prime calculation and the counterexample together give
\[
  m_2+2^{-2^{494}} < \beta_2 \le \frac{907739729}{1073741824}
       <0.845398501494, \qquad m_2=0.8453979555\ldots .
\]
For the explicit lower endpoint, let $I_n$ count the monic binary irreducibles of degree $n$. The factorization of $x^{2^n}-x$ gives $\sum_{d\mid n}d I_d=2^n$. The integer recurrence through $n=480$ gives $R_0+13<2^{494}$ in Section~\ref{sec:tilted}. Moreover
\[
 c_2=\tfrac12\prod_{n\le480}(1+2^{-n})^{-I_n}
 \ge \tfrac12\exp\left(-\sum_{n\le480}1/n\right)>2^{-12},
\]
since $I_n\le2^n/n$. The lower-bound formula for the weighted example proves the displayed inequality.

Thus the gain is strictly positive and smaller than $5.46\times10^{-7}$. The upper endpoint comes from the two-prime obstruction in \emph{Three density questions about geometric progressions in polynomial rings}. We reran that report's exact enumeration through interval length 32. The strict lower inequality is the contribution here. Computing the actual gain to the first decimal at which it differs from $m_2$ remains a specific numerical problem. A smaller degree substitution or a sharper optimized smooth part could improve the constructive lower bound.

The degree bands used in the infinite construction have large gaps. The theorem concerns the limsup of densities in degree balls. Establishing an improvement on a prescribed natural density would require a construction that also controls the intervening cutoffs.

\paragraph{Sources and verification.}
The fixed-radius argument is visible in the readable
\href{https://github.com/lsernau/grh-lean/blob/72fead102f663d5c7a407964a13dc0cf9d977207/grh/src/New/Reduction/Data/Existence.lean\#L80-L116}{source of \texttt{exists\_three\_term\_free}}. The
\href{https://github.com/VladimirReshetnikov/ProveIt/blob/d31039a48d963692e5a5176f47b6603fc9e9d9cf/docs/incoming/grh-progression-densities-and-cubic-prime-race.zip}{earlier ProveIt report} discusses Questions 6.1, 6.2 and 6.4 in the same paper.
\href{https://arxiv.org/pdf/1512.01932}{The original conjecture} is in section 6.
The present disproof is a mathematical argument with finite integer checks. It uses no hypothesis on zeros of $L$-functions. The supplementary analysis explains two transfers of the analytic part of the Lean development.
\end{document}
